\begin{center}
\LARGE
Solver collaborations for non-linear constraints
\end{center}

\begin{center}
\Large
Eric Monfroy

\end{center}

\noindent
Date:  July 19th (Friday)\\
Time:  09:00-09:30\\

\begin{center}
\large
Abstract
\end{center}

Constraint programming is a paradigm based on the notion of constraints and mechanisms for their
resolution. Thus the key point of this class of languages is not only to offer a wide class of
constraints for declarativity reasons, but also to treat them efficiently. For this purpose, the need
for collaboration i.e. combination and cooperation of solvers is widely recognized. This new
concept enables to solve problems that cannot be tackled or efficiently solved with a single solver.
Furthermore the demand for integrating symbolic mathematical tools into automated deduction
system has significantly increased. 

In order to meet these motivations we propose BALI , an environment for designing/executing
solver collaborations. BALI consists of a solver collaboration language and a host language. By
providing several construction primitives (as concurrency, parallelism and sequentiality) and
several control primitives for their composition (as iterator or guarded control), the solver
collaboration language enables to build complex solvers from elementary ones. Then the host
language, which is a constraint programming language, offers several strategies for manipulating
constraints and executing solver collaborations. 

In order to demonstrate the simplicity for designing solver collaborations in BALI, we give
examples over non-linear constraints. The collaborations are made of several solvers as GB (from
J.-C. Faugeere) for Groebner bases, extension of Gaussian elimination for linear equations and
inequations (written in CHRs by T. Fruehwirth) and the symbolic computation software Maple. 

